\documentclass[12pt,a4paper]{article}

\usepackage[T1]{fontenc}
\usepackage{lmodern}
\usepackage{amsmath}

\title{Introduction to LaTeX}
\author{Vikram Mehra}
\date{\today}

\begin{document}

\maketitle

\begin{abstract}
This document demonstrates the basic structure of a LaTeX document.
It introduces document classes, packages, title information,
sections, paragraphs, and mathematical expressions.
\end{abstract}

\section{Introduction}

LaTeX is a document preparation system used for academic writing,
technical reports, research papers, and mathematical documents.

It is especially useful when a document contains equations,
tables, references, and structured sections.

\section{Basic Document Structure}

A LaTeX source file has two main parts: the preamble and the document body.

\subsection{The Preamble}

The preamble appears before the \texttt{\textbackslash begin\{document\}} command.
It contains the document class, packages, and title information.

\subsection{The Document Body}

The document body begins with \texttt{\textbackslash begin\{document\}}
and ends with \texttt{\textbackslash end\{document\}}.
It contains the actual content of the document.

\section{A Mathematical Expression}

LaTeX can display mathematical expressions clearly.

The Pythagorean theorem is:

\[
a^2 + b^2 = c^2
\]

\section{Conclusion}

This example demonstrates the basic structure of a LaTeX document.
More advanced features such as formatting, lists, tables, figures,
and colors are covered in the other practical programs.

\end{document}